\chapter{绪论}

\section{研究背景与意义}

面向技术讨论的在线交流，与日常社交聊天并不一样。在课程项目联调、接口排错、版本确认和开源协作等场景中，参与者反复交换的往往是代码片段、报错日志、接口参数、测试结果和临时结论。消息一多，真正有用的内容就容易被后续讨论淹没，后来加入的人也难以及时接上上下文，只能重新追问，沟通成本随之增加。

现有聊天工具能够承担部分交流任务，但在技术讨论场景中仍有明显不足。通用社交软件偏向日常沟通，对技术内容展示、历史回看和结果留存考虑不多；功能更完整的协作平台覆盖面虽广，却往往带来更高的部署和使用成本。因而，围绕讨论过程组织消息、保留关键结论并支持后续回看，已成为技术交流系统的现实需求\cite{textchat2024}。如果系统还能支持个人电脑或团队内网中的本地部署，就能更好控制消息数据和讨论记录的留存位置。技术讨论中还常有一类高频需求：遇到陌生术语、报错信息或代码片段时，希望直接在当前界面完成划词解释，而不是反复切出聊天窗口搜索。

与此同时，大语言模型和检索增强技术的发展，也让聊天系统有了新的扩展方向。AI 的价值不在于脱离上下文单独回答问题，而在于结合当前会话历史继续参与问答，在需要时完成问题梳理、阶段总结和术语解释\cite{qaoverview2025}。这样一来，用户不必在聊天工具、搜索工具和独立 AI 工具之间来回切换，沟通链路更短，信息也更容易沉淀。因此，设计一套面向技术交流群体的智能聊天室系统，具有较强的实际应用价值和明确的研究意义。

\section{国内外研究现状}

\subsection{即时通讯系统与协作平台的发展}

现有即时通信系统的研究，早已不只停留在“发送一条消息”这一层面。消息留痕、群组组织、文件共享、权限控制和安全通信，已经逐渐成为这类系统的常见能力。针对具体使用场景对聊天系统做定制化裁剪，也是近年来较常见的实现思路。

这对本文的启发比较直接。技术交流群体需要的不是一套笼统的大平台，而是围绕问题讨论本身展开的轻量化工具：消息要清楚，历史要能回看，关键内容要能保留下来，后加入的人也要能尽快看懂前文。

\subsection{Web实时通信技术研究现状}

在 Web 场景下，实时通信问题已经有比较成熟的处理方式。早期轮询和长轮询能够实现近实时效果，但请求频率较高，服务端压力也更明显。WebSocket 普及之后，浏览器与服务器可以保持持续连接，消息推送、在线状态同步和多人广播都更适合放在这一机制下完成\cite{pollingssews2026}。

从工程角度看，当前更常见的做法是把普通请求和实时连接分开使用：前者负责登录鉴权、数据读取和消息落库，后者负责在线状态、增量通知和广播同步。这样处理的好处很务实，链路更清楚，排查问题也更直接。本文的系统设计也是建立在这一思路之上。

\subsection{AI辅助交流技术研究现状}

近年的 AI 研究已经从单纯追问“模型能不能回答问题”，逐渐转向“模型怎样结合外部信息更稳定地回答问题”。已有研究表明，大语言模型在问答、归纳和文本生成方面表现出较强能力，但当输入缺少上下文时，回答仍然容易变得空泛，甚至偏离原问题\cite{gpt4report2023}。

为了解决这一问题，检索增强生成成为近几年较受关注的方向。它的基本思路不是让模型只依赖参数记忆，而是在生成之前补充与当前问题相关的外部信息，从而提高回答的相关性和可用性\cite{ragsurveycn2025}。

针对开发者聊天场景中的阶段总结与多轮追问，已有研究也开始关注摘要方法、对话历史管理以及长上下文利用问题\cite{devchatsummary2025}。

对技术交流场景来说，这一点尤其重要。同一句提问，放在不同聊天室、不同项目阶段里，含义可能完全不同。AI 如果不能接入历史讨论，就很难真正理解当前问题。也正因为如此，本文把 AI 看成聊天链路的一部分，而不是一个脱离场景的独立问答入口。

\section{研究内容与目标}

本文以本人设计并实现的 Mini-Chat-Bar 为对象，围绕技术交流群体的使用特点，完成一套轻量化智能聊天室系统的分析、设计、实现与测试。研究内容主要包括四个方面：一是结合课程项目协作、联调排错和日常技术讨论等场景，梳理目标用户、业务流程与核心需求；二是完成系统总体架构、功能模块和数据组织设计，形成支撑实时通信与 AI 辅助的技术方案；三是实现私聊、聊天室、历史回读、收藏、侧边栏 AI 对话、结合聊天记录的上下文问答、讨论总结与划词解释等功能；四是通过功能测试和性能测试验证系统的可用性与稳定性。

本文的目标不是泛泛讨论“智能聊天”这一概念，而是做出一套能够实际运行的产品原型 Mini-Chat-Bar，并验证它是否适合技术交流场景。该产品既要满足消息收发、历史回看和多人协作等基础交流需要，也要让 AI 以更自然的方式融入讨论过程，例如在侧边栏继续追问、在聊天室里结合历史消息回答问题、对选中的文本做解释，并把有价值的结果保留下来，减少重复沟通。

从价值上看，这套系统可以为课程项目、小型开发团队和技术兴趣社群提供一套门槛相对较低的交流工具，帮助用户更快定位问题、理解术语并回看已有结论。从实现特点上看，本文的侧重点不在于提出新的模型算法，而在于把即时通信、历史消息组织和 AI 辅助放到同一产品链路里：AI 不是独立悬浮在外的问答窗口，而是能够利用聊天记录参与当前讨论，回答结果又能继续被引用、收藏和回看。这种围绕具体使用场景进行整合的做法，也是本课题相对有新意、同时更有现实意义的地方。

\section{论文结构安排}

后文按“需求分析、总体设计、具体实现、测试验证”的顺序展开。

第一章为绪论，交代课题背景，并梳理即时通信、实时通信与 AI 辅助相关的研究现状。

第二章介绍相关技术，说明前端界面组织、后端服务处理、实时通信、数据存储及 AI 检索增强等技术在本项目中的使用方式。

第三章进行系统需求分析，围绕目标用户、典型场景、功能需求与非功能需求，明确系统需要完成和暂不实现的内容。

第四章阐述系统总体设计，包括整体架构、模块划分、数据库设计，以及认证、消息流转、聊天室交互和 AI 辅助等核心流程。

第五章说明系统实现，结合前后端代码与界面，介绍用户、聊天、聊天室、收藏、AI 模块及 Electron 桌面端功能的具体实现。

第六章给出系统测试，从测试环境、功能测试和性能测试三个方面验证运行情况，展示接口、实时通信和 AI 调用链路的测试结果。

第七章为总结与展望，归纳已完成的工作，并说明后续仍可继续改进的方向。

论文正文后的附录统一收录任务书、开题/中期检查、指导记录、审批表与成绩评定等过程性材料，相关内容可结合\hyperref[app:materials]{附录}查阅。
